\chapter{表格数据、星表与 HR 图}

\section{学习目标}

学完本章后，你应该能够：

\begin{itemize}
    \item 理解为什么星表是现代天文学最基础、也最常见的数据形态之一；
    \item 从表观星等、视差和颜色指数出发，计算绝对星等等关键派生量；
    \item 写出构造 HR 图的最小数据处理流程，并解释主序星、巨星和白矮星的大致分布；
    \item 理解 HR 图不仅是“会画一张图”，更是把表格数据转化成恒星物理结构的入口；
    \item 建立对样本筛选、视差质量和选择效应的初步敏感性。
\end{itemize}

进入本章前，建议你已经完成 Part 0 的表格读取与基础控制流训练，至少能顺手写出 \texttt{Path}、\texttt{csv.DictReader}、简单筛选和派生列计算。本章会把这些最小能力直接用在真实星表上。

\section{为什么星表是天文数据分析的起点}

在很多天文项目中，真正先进入你手里的并不是图像，而是一张表。它可能来自：

\begin{itemize}
    \item Gaia 的恒星测光与视差目录；
    \item 某次交叉匹配后的目标样本；
    \item 某个观测计划整理出的候选源列表；
    \item 经过前处理的教学用 CSV 文件。
\end{itemize}

和图像、光谱相比，星表更适合初学者进入真实数据工作流，因为它已经把观测结果组织成“按行代表对象、按列代表属性”的结构化形式。只要你能读懂这种结构，后面的筛选、统计、可视化、机器学习和案例分析都会自然接上。

\begin{defbox}[星表]
星表是一种以对象为行、以观测属性或派生属性为列的数据组织形式。它是现代天文学中最基础的数据载体之一。
\end{defbox}

\section{星表中的最关键字段：不是每一列都一样重要}

一个真实星表可能有几十甚至上百列字段，但教学里更重要的不是“一次全看完”，而是先识别哪些字段真正支撑当前物理问题。

在本章配套数据 \nolinkurl{stellar_photometry_demo.csv} 中，我们重点关注：

\begin{itemize}
    \item \texttt{source\_id}：对象标识；
    \item \texttt{bp\_rp}：颜色指数，近似反映恒星表面温度；
    \item \texttt{parallax\_mas}：视差，帮助估计距离；
    \item \texttt{g\_mag}：表观星等；
    \item \texttt{stellar\_class}：教学标签，用于帮助解释分布。
\end{itemize}

这里最值得强调的是：\textbf{星表分析并不是“看到很多列就都拿来用”，而是要先明确哪几列能回答当前科学问题。}

\subsection{星表字段契约：先确认列能否支撑问题}

进入任何星表分析前，都应该先写下一个很小的字段契约。所谓字段契约，就是明确本章分析需要哪些列、这些列的单位是什么、允许的数值范围是什么，以及缺失或异常时如何处理。

对本章 HR 图任务来说，最小契约可以写成：

\begin{itemize}
    \item \texttt{bp\_rp} 必须是有限数值，单位为 mag；
    \item \texttt{g\_mag} 必须是有限数值，单位为 mag；
    \item \texttt{parallax\_mas} 必须为正，单位为 mas；
    \item 若存在 \texttt{parallax\_error\_mas}，应优先检查视差信噪比；
    \item 派生出的 \texttt{absolute\_g\_mag} 不能和原始表观星等混用。
\end{itemize}

\begin{examplebox}[最小字段契约检查]
\begin{lstlisting}[style=pythonstyle]
required_fields = ["bp_rp", "g_mag", "parallax_mas"]
missing = [name for name in required_fields if name not in rows[0]]
if missing:
    raise ValueError(f"Missing required fields: {missing}")
\end{lstlisting}
\end{examplebox}

这个检查看起来朴素，却能提前挡住很多后续错误。越是进入真实 Gaia 表，字段越多，越需要先把“当前分析真正依赖什么”说清楚。

\section{从表观星等到绝对星等：为什么必须做这个变换}

如果你直接把颜色指数和表观星等画在一起，图形当然也可能有一定结构，但它混入了距离效应。天体看起来暗，未必本征上暗；也可能只是更远。

因此，构造 HR 图时更常用的是绝对星等 $M$。其定义来自距离模数公式：

\[
M = m - 5\log_{10}\left(\frac{d}{10\,\mathrm{pc}}\right),
\]

其中 $m$ 是表观星等，$d$ 是距离（单位：pc）。如果进一步用上一章学过的视差关系

\[
d(\mathrm{pc}) = \frac{1000}{p(\mathrm{mas})},
\]

代入上式，就得到：

\[
M = m - 10 + 5\log_{10} p(\mathrm{mas}).
\]

这就是配套 notebook 中最关键的派生量计算之一。

\begin{examplebox}[从星表字段计算绝对星等]
\begin{lstlisting}[style=pythonstyle]
import csv
import math
from pathlib import Path

data_path = Path("../../data/small/stellar_photometry_demo.csv")
with data_path.open(newline="", encoding="utf-8") as f:
    rows = list(csv.DictReader(f))

for row in rows:
    row["bp_rp"] = float(row["bp_rp"])
    row["parallax_mas"] = float(row["parallax_mas"])
    row["g_mag"] = float(row["g_mag"])
    row["absolute_g_mag"] = (
        row["g_mag"] - 10.0 + 5.0 * math.log10(row["parallax_mas"])
    )
\end{lstlisting}
\end{examplebox}

这一小步本身就非常教材化，因为它把三个层次连起来了：

\begin{itemize}
    \item 观测量：表观星等与视差；
    \item 物理目标：本征亮度的比较；
    \item 代码实现：在表格上生成派生列。
\end{itemize}

\section{绝对星等的不确定度与视差质量}

绝对星等并不是一个“算出来就完了”的派生量。只要视差有误差，它就会跟着带上不确定度。由

\[
M = m - 10 + 5\log_{10} p(\mathrm{mas})
\]

可得，在忽略相关项的情况下，

\[
\sigma_M^2 \approx \sigma_m^2 + \left(\frac{5}{\ln 10}\frac{\sigma_p}{p}\right)^2.
\]

这说明两件事：

\begin{itemize}
    \item 视差越小，绝对星等越敏感；
    \item 如果视差信噪比很低，直接把 $1000/p$ 当作距离会变得非常不稳定。
\end{itemize}

\begin{warnbox}[负视差和低信噪比视差]
在真实 Gaia 数据里，视差可能接近零，甚至出现负值。这并不表示“天体真的在负距离”，而是说明测量噪声已经大到不能直接用简单反演。教学数据里可以先做正视差筛选，但真实工作通常还需要更谨慎的距离估计方法。
\end{warnbox}

\begin{protip}[最小质量筛选]
教学 notebook 中最常见的最小筛选条件可以先写成：数值有限、视差为正、视差信噪比足够高。真实项目里当然还要看更多质量标志，但这三步已经能避免很多明显错误。
\end{protip}

\subsection{最小筛选代码}

本章的教学样本没有把 Gaia 全部质量标志都展开出来，所以最小筛选可以先写成“关键字段存在、视差为正、派生量可计算”。真实数据里通常还要继续看视差信噪比、天体测量拟合质量和消光影响，但这里先把最小流程跑通。

\begin{examplebox}[HR 图前的最小筛选]
\begin{lstlisting}[style=pythonstyle]
import math

selected_rows = [
    row for row in rows
    if row["parallax_mas"] > 0.0
    and math.isfinite(row["bp_rp"])
    and math.isfinite(row["g_mag"])
]
\end{lstlisting}
\end{examplebox}

\section{颜色指数为什么可以放在横轴}

在严格恒星物理里，HR 图常把温度或光谱型放在横轴；但在现代巡天数据处理里，更常见的是用颜色指数（例如 \texttt{BP-RP}）作为温度 proxy。

颜色指数的直觉是：

\[
\mathrm{color} = m_{\mathrm{blue}} - m_{\mathrm{red}}.
\]

若一个恒星在蓝端更亮，那么它的颜色指数通常更小，看起来更“蓝”；若红端更亮，则颜色指数更大，看起来更“红”。在许多教学数据和初步分析里，这足以给出温度排序的第一近似。

因此，HR 图在教学版中通常可以理解为：

\begin{itemize}
    \item 横轴：颜色指数，近似代表温度；
    \item 纵轴：绝对星等，近似代表本征亮度。
\end{itemize}

\section{HR 图到底在告诉我们什么}

HR 图最经典的教学价值，在于它把原本分散在多列字段里的信息，压缩到一张物理结构图里。

在这个平面上，常见的几个区域包括：

\begin{itemize}
    \item \textbf{主序星}：从蓝亮到红暗形成一条主带，是恒星生命周期中最常见的阶段；
    \item \textbf{巨星}：温度未必极高，但本征光度明显更大，因此位于主序带上方；
    \item \textbf{白矮星}：通常较蓝但很暗，因此落在图的另一片区域。
\end{itemize}

即使在很小的教学样本中，你也可以开始形成这些物理直觉；当样本扩大到真实 Gaia 数据时，这些点会变成非常清晰的统计结构。

\section{为什么图上的主序带会有厚度}

初学者常常会问：如果 HR 图描述的是恒星物理规律，为什么主序带不是一条细线？答案是，它本来就不该是一条细线。图上的厚度来自多个来源：

\begin{itemize}
    \item 测量误差会让点在横纵方向上扩散；
    \item 未解析双星会让系统看起来更亮；
    \item 金属丰度差异会改变颜色和亮度关系；
    \item 星际消光会让天体显得更暗、更红；
    \item 样本选择和观测极限会改变图中点的密度分布。
\end{itemize}

\begin{defbox}[消光]
消光是星际尘埃对光的吸收和散射效应。它通常会让天体看起来更暗，也更偏红，因此在 HR 图里往往会把点沿着某个斜方向推移。
\end{defbox}

这就是为什么 HR 图的解读不能只停留在“哪一类点在什么位置”，还要继续问：这些位置变化里，有多少来自真实恒星性质，有多少来自观测条件。

\section{为什么纵轴常常要反向显示}

绝对星等的定义对初学者有一个经典“反直觉”之处：数值越小，天体越亮。因此在许多 HR 图中，我们会把纵轴反向显示，让亮星出现在图的上方。这并不是美观上的约定，而是为了让图形结构更符合人类对“亮在上、暗在下”的阅读习惯。

如果不用反向纵轴，图仍然是数学上正确的，但物理直觉会更难形成。

\begin{examplebox}[教学版 HR 图代码]
\begin{lstlisting}[style=pythonstyle]
ax.scatter(color_values, absolute_magnitudes, color="tab:orange")
ax.set_title("Teaching HR Diagram Demo")
ax.set_xlabel("BP - RP [mag]")
ax.set_ylabel("Absolute G magnitude [mag]")
ax.invert_yaxis()
\end{lstlisting}
\end{examplebox}

\section{从表格到图像：一个最小算法流程}

本章真正想让学生掌握的，并不是某一条 Matplotlib 命令，而是构造 HR 图的最小流程：

\begin{enumerate}
    \item 读入星表；
    \item 把关键列转成数值；
    \item 检查视差是否为正、是否数量级合理；
    \item 计算绝对星等；
    \item 按需要筛掉不可信样本；
    \item 再把颜色指数和绝对星等画出来；
    \item 最后回到物理解释：这些点为什么会落在这些区域。
\end{enumerate}

这条算法流程，在形式上和后面的机器学习 pipeline 并没有本质区别：先清楚数据，再构造派生量，再做可解释输出。

\section{HR 图解释的三个层次}

画出 HR 图以后，解释也应分层进行，而不是直接跳到“这是主序星”。一个更稳妥的阅读顺序是：

\begin{enumerate}
    \item \textbf{数据层}：点是否来自可信视差、有限颜色和合理星等；
    \item \textbf{图形层}：点云是否形成主带、上方亮星区或蓝暗区域；
    \item \textbf{物理层}：这些区域是否可以解释为主序星、巨星或白矮星，并且是否受到消光、双星或选择效应影响。
\end{enumerate}

这三个层次不能互相替代。数据层没有过关时，图形层的结构可能只是错误计算的结果；图形层没有描述清楚时，物理层的解释就容易变成背诵；物理层不讨论限制时，HR 图又会退化成一张漂亮散点图。

\begin{protip}[先描述，再解释，最后限定]
一段好的 HR 图解读通常先说明轴和样本筛选，再描述点云结构，然后给出物理解释，最后补一句当前解释的限制。这个顺序比直接说“这里是主序星”更接近科研写作。
\end{protip}

\section{样本筛选与选择效应：为什么“画出来”还不够}

在真实数据里，HR 图的形状还会受到选择效应影响。例如：

\begin{itemize}
    \item 视差误差大的远处样本可能让绝对星等估计变得不稳定；
    \item 亮源和暗源的探测完整性不一样；
    \item 某些波段下的颜色测量对尘埃消光敏感。
\end{itemize}

如果你把这些效应忽略掉，很容易把观测偏差误读成恒星物理差异。特别是在视差质量较差的区域里，直接用 $1/p$ 得到的距离会放大噪声，最后让整张 HR 图看起来“散”，但这种散并不一定是天体本身真的更复杂。

本章的教学数据刻意保持简单，但教学上必须明确：\textbf{图形中的结构不只来自恒星物理，也可能来自观测条件和样本筛选。}

\section{配套 notebook 做了什么}

配套 notebook 采用一条清楚的教学主线：

\begin{itemize}
    \item 读取一个小型 Gaia 风格星表；
    \item 生成绝对星等派生列；
    \item 检查最亮、最暗、最蓝样本的基本性质；
    \item 在安装了 \texttt{matplotlib} 的环境中画教学版 HR 图；
    \item 最后结合标签做物理解释，而不是只停留在“图画出来了”。
\end{itemize}

\section{常见错误}

\begin{warnbox}[星表与 HR 图中的典型问题]
\begin{itemize}
    \item 把表观星等直接当成绝对星等；
    \item 忘记视差单位是毫角秒；
    \item 看到散点图后只会描述颜色和形状，不会回到恒星物理解释；
    \item 完全忽略视差质量和样本筛选问题；
    \item 在很小的教学样本上过度解读微小分布差异。
\end{itemize}
\end{warnbox}

\section{AI 助手如何帮助你}

在这一章，AI 助手适合帮助你：

\begin{itemize}
    \item 帮你把星表读入与派生量计算代码整理得更清楚；
    \item 检查某个公式中视差和距离的单位是否一致；
    \item 提醒你 HR 图中是否缺少轴标签、单位或纵轴反转；
    \item 帮你把“主序星/巨星/白矮星”这些区域写成更清楚的图像解读说明。
\end{itemize}

但图中出现的结构到底意味着什么、某个样本是否应该被筛掉、以及一个区域的物理解读是否合理，仍然必须由你自己做判断。

\section{练习}

\begin{exercisebox}[基础练习]
把样本按 \texttt{stellar\_class} 分组，并用一段短文字比较三类恒星在颜色和绝对星等上的差异。不要只说“不同”，而要说明差异大致体现在哪个方向。
\end{exercisebox}

\begin{exercisebox}[进阶练习]
从公式
\[
M = m - 5\log_{10}(d/10)
\]
出发，重新推导出
\[
M = m - 10 + 5\log_{10} p(\mathrm{mas})
\]
并说明其中为什么会出现常数项 $10$。
\end{exercisebox}

\begin{exercisebox}[开放练习]
假设你手里有一个更真实的 Gaia 子样本。请写一个“画 HR 图前的最小检查清单”，至少包含：视差质量、字段单位、缺失值、派生量公式和选择效应提醒。
\end{exercisebox}

\section{本章小结}

HR 图不是一张“天文经典图片”，而是一条完整的数据到物理的工作流：从星表字段出发，通过距离修正构造绝对星等，再把颜色和本征亮度组织成恒星演化结构。只要你真正理解这条链条，后面的案例分析和机器学习章节就会更有物理根基。
